\subsubsection{7.4 Was nicht automatisch
folgt}\label{was-nicht-automatisch-folgt}

\begin{itemize}
\tightlist
\item
  \textbf{Zeit:} Für einen Grundzustand ist
  \codeword{e\textasciicircum{}\{-itH\}Ω=e\textasciicircum{}\{-itE0\}Ω};
  sein physischer Zustand bleibt gleich. Das ist kein eigener Zeitpfeil.
  Nichtstationäre Präparation, Korrelationen und Records müssen
  ausgewiesen werden.
\item
  \textbf{Eichfeld:} Für bloße Basiswechsel \codeword{Uxy=Vx†Vy}
  teleskopiert jede geschlossene Holonomie zu I. Echte Krümmung braucht
  zusätzliche Verbindungs-/Transportdaten oder einen nachgewiesenen
  Mechanismus.
\item
  \textbf{Spin:} Zwei Orientierungsmarken liefern nicht von selbst einen
  SL(2,C)-Spinor, dessen Transformationsgesetz und Kinetik. Der
  zusätzliche antisymmetrische Zweierfaktor muss unabhängig konstruiert
  sein.
\item
  \textbf{Raum:} Minimale Operationskosten sind ohne Reversibilität
  zunächst gerichtete Kosten, nicht notwendig eine symmetrische Metrik.
  Kubisches Volumenwachstum allein beweist keine 3D-Mannigfaltigkeit,
  keine Lorentzstruktur und keinen gemeinsamen Lichtkegel.
\item
  \textbf{Gravitation:} Dynamische Beziehungen sind ein möglicher
  Träger, aber ein masseloser Spin-2-Sektor, zwei Helizitäten,
  Energiepositivität und universelle konsistente Kopplung bleiben zu
  beweisen.
\item
  \textbf{Arithmetik:} Primitive Schleifen besitzen generisch neue
  gemischte primitive Zyklen. Sie sind nicht automatisch Primzahlen. Die
  Xi-Determinante bleibt eine unbewiesene vollständige
  Operatoridentität. Der getrennte RH-Anhang dokumentiert Vorarbeiten,
  Beispiele und Quellenstatus.
\end{itemize}

Der Text wird deshalb als \textbf{präzisierter globaler
Forschungsansatz} integriert, nicht als Nachweis, dass derselbe
unbekannte Baustein schon alle T1-T8-, RH-, Faktorisierungs-, P/NP- und
Hylæan-Fragen beantwortet.

\subsection{8. T1-T8: gemeinsamer Fortschritt ohne
Statusinflation}\label{t1-t8-gemeinsamer-fortschritt-ohne-statusinflation}

{\def\LTcaptype{none} % do not increment counter
\begin{longtable}[]{@{}
  >{\raggedright\arraybackslash}p{(\linewidth - 4\tabcolsep) * \real{0.3333}}
  >{\raggedright\arraybackslash}p{(\linewidth - 4\tabcolsep) * \real{0.3333}}
  >{\raggedright\arraybackslash}p{(\linewidth - 4\tabcolsep) * \real{0.3333}}@{}}
\toprule\noalign{}
\begin{minipage}[b]{\linewidth}\raggedright
Front
\end{minipage} & \begin{minipage}[b]{\linewidth}\raggedright
In dieser Revision weitergekommen
\end{minipage} & \begin{minipage}[b]{\linewidth}\raggedright
Noch erforderlicher Nachweis
\end{minipage} \\
\midrule\noalign{}
\endhead
\bottomrule\noalign{}
\endlastfoot
T1 - Ursprung und Auswahl & Explizite Identität zwischen Quell-Clock und
innerer Spinwirkung; Operationsverträge genauer & Primitive Auswahl von
P1/P2, Träger, erlaubten Instrumenten und Zustand \\
T2 - Half-Charge/E8-Feld & Bestehender geladener Pol bleibt abgesichert,
innere Darstellung präzisiert & Renormiertes Half-Charge-Feld mit
Energie-/Adjungiertenkontrolle und E8-/Clock-Adapter \\
T3 - gemeinsamer 3+1D-Parent & Kleiner exakter Paartransferzeuge und
korrigierter Überlappungstest & Quellenseitig ausgewählter räumlicher
Parent auf einem gemeinsamen Hilbertraum und Zustand \\
T4 - chirale Materie & Basisunabhängiger Ausschluss eines konkreten
Ein-Kopie-Vektoradapters & Konsistenter chiraler Feldadapter,
Maß/Anomalien/Index und Spiegelentkopplung \\
T5 - Kontinuum und Dynamik & Exakte endliche Clock-/Casimir-Kombination
und bedingter invarianten Transferblock & Kontrollierter
wechselwirkender Grenzübergang, Lorentzverhalten, Clustering und
Streuung \\
T6 - Kopplungen und Texturen & Keine neue Herleitung; neue
Mischparameter sichtbar zusätzlich & Herkunft aller Kopplungen und
vollständiger Neutrinostruktur/-skala \\
T7 - Gravitation & Der falsche pauschale Tensor-Ausschluss bleibt
zurückgenommen & Masseloser dynamischer Spin zwei, beide Helizitäten,
konsistente universelle Kopplung im selben Parent \\
T8 - Präparation und Auslesung & Gram-/Antworttest gemeinsam formuliert;
spezielle Blockerressource konkret & Natives geladenes Instrument,
Quellzustandswahl, Records und ein gemeinsames Auslesefunktional \\
\end{longtable}
}

Keiner der acht vollständigen Abschlussverträge wird in dieser Revision
auf „gelöst`` gesetzt. Gelöst wurden ausdrücklich benannte Teilfragen.

\subsection{9. Die nächsten drei Arbeiten - klein, gemeinsam und
entscheidbar}\label{die-nuxe4chsten-drei-arbeiten---klein-gemeinsam-und-entscheidbar}

\textbf{A. Den gemeinsamen Quellenvertrag wirklich hinschreiben.} Zwei
lokale Einbettungen, ihr CAR-Gram, globale Parität und Ladung,
ursprüngliche Operationen und ein einziges H angeben. Abbruch für eine
Kandidatenfassung, wenn sie nur zwei Koordinatenansichten desselben
unveränderten Signals umbenennt. Die gemeinsamen Generatoren müssen
einen nachweislich neuen Antwortprozess ermöglichen, nicht nur eine
andere Beschriftung.

\textbf{B. Die fehlende Zwischenoperation ausweisen.} Der
Vierzustandszeuge gibt eine konkrete positive Zieloperation vor. Zu
prüfen ist, ob der Compiler ein geeignetes Mischinstrument oder eine
andere geteilte Paarstruktur wirklich liefert. Clock und Casimire
alleine sind kein Herkunftsbeweis. Zeigt ein erhaltener gemeinsamer
Generator die Unmöglichkeit für einen festen Operationssatz, endet
dessen weitere Kommutator-Suche; dann muss ein anderer bereits
vorhandener Quellbestandteil benannt werden.

\textbf{C. Dieselbe Anregung und denselben Zustand behalten.} Auf dem
Kandidaten \codeword{G\_AB(0)} und \codeword{G\_AB(t)} einschließlich
vollständiger Restantwort bestimmen. Der N=9-Zeuge darf nicht still zum
N=64-Lochpol umgedeutet werden. Ein minimaler Lorentz-/Kinetikadapter
muss die nachgerechneten Tensorbedingungen erfüllen. Erst mit diesem
gemeinsamen Anschluss trägt eine großräumige Skalierungsrechnung.

Der stärkste Vereinfachungsgewinn dieser Runde ist damit konkret:
\textbf{Clock und innere Symmetrie sind nicht zwei voneinander
unabhängige Mechanismen; reine Überlappung und echte Bewegung sind
dagegen zwei verschiedene Dinge.} Die weitere Suche sollte diese erste
Identität nutzen und die zweite Unterscheidung im selben kleinen Modell
erzwingen.

\subsection{10. Belege und
Reproduzierbarkeit}\label{belege-und-reproduzierbarkeit}

Die drei Eingangstexte bleiben unverändert im Prüfpaket. Ein Teil-Audit
reproduziert die 1127 gelieferten Guards; eigene Programme behandeln
Clock/Casimir, Quellenreichweite, direkte Krylov-Kontraktion,
basisunabhängige Feldbedingungen, den tatsächlichen Vierzustandskanal
und die Überlappungsgegenkontrollen. Exakte und numerische Teile sind in
den jeweiligen Ergebnisdateien getrennt markiert. Die abschließenden
Zählungen und Bytevergleiche stehen im Auslieferungs-/Replayprotokoll;
interne erneute Aufrufe der 944 Guards werden nicht mehrfach als neue
Tests gezählt.

Die drei ausführlichen parallelen Berichte und der
Überlappungs-/RH-Nachtrag folgen im vollständigen Dokument. Ihr Status
ersetzt widersprechende historische Deutungen, nicht unveränderte ältere
Beweise. Das native frühere Archiv bleibt zusätzlich unverändert im
Paket; dessen große frühere Enumeration wird nicht als in dieser
Revision erneut durchgeführt ausgegeben.

\begin{center}\rule{0.5\linewidth}{0.5pt}\end{center}

\section{Anhang A: Quellenprüfung Symmetrie und
Verfügbarkeit}\label{anhang-a-quellenpruxfcfung-symmetrie-und-verfuxfcgbarkeit}

\section{Quellenprüfung: Quellsymmetrie, Verfügbarkeit und
Lorentztypen}\label{quellenpruxfcfung-quellsymmetrie-verfuxfcgbarkeit-und-lorentztypen}

Stand: 15. September 2026. Eigenständiger Teil-Audit der ersten neuen
Anlage und dreier Programme. Keine Änderung der gelieferten Programme,
Ergebnisdateien oder Dokumente. Keine Aussage, dass T1-T8 geschlossen
seien.

\subsection{1. Reproduktion}\label{reproduktion-1}

Alle drei Programme wurden vollständig gelesen und anschließend in
isolierten Kopien ausgeführt. Die relevanten Tensoren wurden an ihre von
den unveränderten Programmen erwarteten relativen Orte kopiert. Normaler
und optimierter Lauf reproduzieren jeweils die gelieferten
Ergebnisdateien byteidentisch, ohne Standardfehlerausgabe:

{\def\LTcaptype{none} % do not increment counter
\begin{longtable}[]{@{}lrl@{}}
\toprule\noalign{}
Programm & Gelieferte Prüfbedingungen & Reproduktion \\
\midrule\noalign{}
\endhead
\bottomrule\noalign{}
\endlastfoot
\codeword{operation\_symmetry.py} & 944 & PASS, identisch \\
\codeword{symmetry\_availability.py} & 31 & PASS, identisch \\
\codeword{lorentz\_types.py} & 152 & PASS, identisch \\
Summe & 1127 & normal und \codeword{-OO} \\
\end{longtable}
}

Die Zusatzprüfung \codeword{check\_scope.py} rechnet kleine exakte
Gegenprüfungen zur Reichweite der Aussagen. Das ausführbare
Reproduktionsprogramm ist \codeword{replay.py}; Quellenpins, Kopien,
Originalberichte und Ausgaben liegen in diesem Auditordner.
\codeword{replay\_receipt.json} dokumentiert alle Pins und die
Unverändertheit der Eingaben.

Programm-Pins:

{\def\LTcaptype{none} % do not increment counter
\begin{longtable}[]{@{}
  >{\raggedright\arraybackslash}p{(\linewidth - 2\tabcolsep) * \real{0.5000}}
  >{\raggedright\arraybackslash}p{(\linewidth - 2\tabcolsep) * \real{0.5000}}@{}}
\toprule\noalign{}
\begin{minipage}[b]{\linewidth}\raggedright
Quelle
\end{minipage} & \begin{minipage}[b]{\linewidth}\raggedright
SHA-256
\end{minipage} \\
\midrule\noalign{}
\endhead
\bottomrule\noalign{}
\endlastfoot
\codeword{operation\_symmetry.py} &
\codeword{ec93f759d6582f42542a597bf916ee2a62a51c510240aa777ae439e90fe28382} \\
\codeword{symmetry\_availability.py} &
\codeword{88cc3e3361a198fe7b1b32a7174f3dd45f341106978d8abd8d5722e0cb631967} \\
\codeword{lorentz\_types.py} &
\codeword{43777c81ea8ba2d2ae12b711321a54abd7f7152a326611a296598d5c5dd463da} \\
Nativer Tensor &
\codeword{3f00a0892157a691e4ef26920c3702772c7bf76a250fad98ff04a079bb64b763} \\
\end{longtable}
}

Die Erzählung über zwei unabhängig entstandene, vollständig
übereinstimmende Programme ist durch diese Reproduktion \textbf{nicht}
geprüft: Vorgelegen hat die erhaltene Fassung mit 944 Guards, nicht die
vorherige Fassung mit 528 Guards.

\subsection{2. Bestätigte Mathematik}\label{bestuxe4tigte-mathematik}

Der N=3-Raum zerfällt als Darstellung von Spin(10) × SU(4) in sieben
Isotypien. Drei haben Multiplizität zwei, die anderen vier Multiplizität
eins. Die zwei Fockteile Λ³(C⁶⁴) und C⁶⁰⊗C⁶⁴ sind jeweils
multiplizitätsfrei; der gesamte N=3-Raum ist es ausdrücklich
\textbf{nicht}.

Die irreduziblen Dimensionen sind 24000, 11200, 2688, 2880, 576, 320,
64; die mit Multiplizitäten gewichtete Summe ist 45504. Die
Spin(10)-Darstellung mit Dimension 16 in BF und Λ³ ist in den
angegebenen Gewichten die konjugierte Spinordarstellung; die
Dimensionskurzschreibweise \codeword{16} darf diese Information nicht
ersetzen.

Bestätigt sind:

\begin{itemize}
\tightlist
\item
  die exakte Zerlegung über Gewichtsmultiplizitäten und
  Racah-Alternation;
\item
  \codeword{C(560)=117/4} in der erklärten Casimir-Normierung;
\item
  die Quellenintertwiner für alle 60 Lie-Erzeuger;
\item
  \codeword{rank C3=3776}, \codeword{dim ker C3=37888};
\item
  Spektrum von \codeword{C3 C3†}: 0,7,10,12 mit Multiplizitäten
  64,2880,576,320;
\item
  die drei dunklen Typen mit Gesamtcasimir 45, aber getrennten
  Spin-Casimiren 141/4,117/4,165/4;
\item
  für den \textbf{gewährten} Satz X,Nb eine assoziative Algebra der
  komplexen Dimension 14;
\item
  für alle punktweise symmetrieinvarianten Operatoren die Dimension 16;
\item
  nach zusätzlicher Gewährung aller 60 Gruppenerzeuger die assoziative
  Dimension 743583744 und Kommutantdimension 7.
\end{itemize}

Die Eigenwertzuordnung im Code prüft pro S-Eigenwert einen Vektor, nicht
eine vollständige Basis dieses Eigenraums. Zusammen mit der unabhängig
geprüften Multiplizitätsfreiheit, der Kommutation mit allen Erzeugern
und der eindeutigen Dimensionszuordnung der BF-Teilsummen ist die
Zuordnung trotzdem abgesichert. Die Guard-Beschriftung »entire
eigenspace« ist als direkte Beschreibung dieser einen Rechnung zu weit.

\subsection{3. Invarianz und Kovarianz
auseinanderhalten}\label{invarianz-und-kovarianz-auseinanderhalten}

Mit \codeword{A0=Alg*(X,Nb)} und Gruppenwirkung ρ bezeichne

\codeword{C = End\_G(H3) = ⊕\allowbreak{}\_i End(C\textasciicircum{}\{m\_i\}) ⊗ I\_\{d\_i\}}.

Dann gilt \codeword{A0 ⊂ C}, \codeword{dim A0=14} und
\codeword{dim C=Σ m\_i²=16}. Jedes Wort aus X und Nb kommutiert mit
ρ(G). Nichtzentrale Lie-Erzeuger können deshalb aus diesem Alphabet
nicht entstehen. Dieser Ausschluss ist korrekt und bleibt richtig, wenn
nur ein fixes H statt zweier unabhängiger Kontrollen gewährt ist.

Nach Gewährung der Gruppenerzeuger entsteht dagegen

\codeword{Afull = ⊕\allowbreak{}\_i End(C\textasciicircum{}\{m\_i\} ⊗ V\_i)}.

Diese Algebra ist \textbf{unter Gruppenkonjugation stabil}, aber ihre
Operatoren sind nicht sämtlich invariant. In
\codeword{symmetry\_availability.py} ist die Kennzeichnung
\codeword{symmetry\_invariant: true} für diese Zeile daher falsch. Der
Kommutant von Afull ist die sieben-dimensionale skalare Blockmitte.

Sieben ist ein Boden, solange alle zusätzlichen Operationen diese sieben
Isotypie-Projektoren erhalten; insbesondere senken zusätzliche
punktweise G-invariante Operationen den bereits erreichten Kommutanten
nicht weiter. Sieben ist aber \textbf{kein} universeller Boden für unter
G stabile Algebren oder alle denkbaren nativen Verträge.
\codeword{End(H3)} ist unter jeder Gruppenkonjugation stabil und hat nur
den skalaren Kommutanten. Ein exaktes 2×2-Gegenmodell in der
Zusatzprüfung macht die Unterscheidung direkt sichtbar.

Die Zahlen sind zudem Dimensionen \textbf{komplexer assoziativer
Sternalgebren}, keine nachgewiesenen Dimensionen einer dynamischen
Lie-Algebra oder erreichbaren Unitärgruppe. Burnside liefert nicht
automatisch beliebige ausführbare Gatter auf jedem irreduziblen Block.

\subsection{4. Verfügbarkeit ist keine universelle
Ja/Nein-Frage}\label{verfuxfcgbarkeit-ist-keine-universelle-janein-frage}

Ein vorgegebener Hamiltonoperator allein gewährt nicht bereits das
unabhängige Schalten von X und Nb. Am erklärten Punkt g/Δ=1/20 besitzt
ein fixes H im N=3-Raum acht verschiedene Energien, also nur eine
acht-dimensionale kommutative Spektralalgebra. Die großzügigere Annahme
unabhängiger Kontrollen X,Nb ergibt 14.

Eine treue, auf die Multiplizitätsräume reduzierte rationale
Matrixrechnung ergibt:

{\def\LTcaptype{none} % do not increment counter
\begin{longtable}[]{@{}lr@{}}
\toprule\noalign{}
Zusätzlich gewährter Satz & Assoziative Dimension \\
\midrule\noalign{}
\endhead
\bottomrule\noalign{}
\endlastfoot
Nur fixes H=Nb+X/20 & 8 \\
X,Nb & 14 \\
X,Nb plus Projektor auf genau einen dunklen Typ & 15 \\
X,Nb plus Spin(10)-Casimir & 16 \\
X,Nb plus SU(4)-Casimir & 16 \\
X,Nb plus Gesamtcasimir & 14 \\
\end{longtable}
}

Bereits \textbf{einer} der getrennten Casimire reicht zur vollständigen
Trennung der drei dunklen Typen. Die zwei fehlenden Algebradimensionen
sind kein Beweis, dass genau zwei physische Griffe fehlen. Die Dimension
15 zeigt eine mögliche invariante Zwischenstufe. Die behauptete
universelle Dichotomie »getrennte Griffe ja oder unverändert 14« und das
\codeword{IF AND ONLY IF} zur Verfügbarkeit der vollen Symmetrie sind
daher zu stark.

Die zusätzlichen Projektoren und Casimirkontrollen werden hier nur als
algebraische Gegenbeispiele verwendet. Ihre Implementierung aus dem
Compiler wird nicht behauptet.

\subsection{5. Der Clock-Satz ist im gelieferten Checker nicht
geprüft}\label{der-clock-satz-ist-im-gelieferten-checker-nicht-gepruxfcft}

\codeword{operation\_symmetry.py:478} schreibt den Normalisator- und
Nichts-hinzufügen-Satz als \codeword{need(True, ...)}. Insgesamt enthält
diese Datei neun solche als Prüfbedingungen gezählten theoretischen
Folgerungen. Einige sind durch die vorausgehende Mathematik gut
begründet; der bloße grüne Guard zertifiziert sie jedoch nicht
unabhängig.

Die direkte Clock-Definition liegt in
\codeword{universalraum-native-operations-ground-response-20260915/common.py},
Funktion \codeword{clock\_lift}, ab Zeile 187: eine Permutation der fünf
Oszillatorachsen wird mit Exterior-Vorzeichen auf den geraden Spinorraum
gehoben; \codeword{GF=G16⊗I4}. Der Bosonlift wird entsprechend gebaut
und die Tensorintertwining-Gleichung geprüft. Eine neue direkte Prüfung
der Clock-/Casimir-Beziehungen erfolgt im parallelen Hauptstrang, nicht
in diesem Audit.

